% !PRINTED 402
\weberpage{402}
%<<CONT
\[
x_1 = a_1 + \frac{1}{x_2}
\]
in $f_2(x_2)$. The equation $f_2 = 0$ then again has at least
one root that is greater than 1, and which therefore lies between
the two positive integers $a_2$ and $a_2 + 1$, and so on.
Finally one will in this way come to isolate a single
root, and one obtains in the sequence of convergents
\[
(a_0), \quad (a_0, a_1), \quad (a_0, a_1, a_2), \quad (a_0, a_1, a_2, a_3) \ldots
\]
a sequence of rational numbers which lie alternately above and below
one of the roots $x$ and approach this root without limit.
If one has advanced in this way to the convergent
\[
(a_0, a_1, a_2 \ldots a_{\nu-1}) = \frac{P_\nu}{Q_\nu}
\]
then this differs from the true value
of $x$ by less than $1 : Q_\nu^2$. The size of the denominator thus gives
directly a measure of the accuracy already attained.

In the ordinary cases one will obtain the numbers $a_0, a_1, a_2 \ldots$
simply by attending to the changes of sign of the
functions $f, f_1, f_2 \ldots$

To test this procedure on an example, we take
the cubic equation which we have already considered
earlier:
\[
x^3 - 2x - 2 = 0.
\]
It has a root between 1 and 2, so we have to set $a_0 = 1$.

If one actually carries out the transformation indicated above,
one obtains the following chain of equations:
\[
\begin{aligned}
x^3 - 2x - 2 &= 0, & a_0 &= 1\\
3x^3 - x^2 - 3x - 1 &= 0, & a_1 &= 1\\
2x^3 - 4x^2 - 8x - 3 &= 0, & a_2 &= 3\\
9x^3 - 22x^2 - 14x - 2 &= 0, & a_4 &= 1\\ % [sic] printed order a_4 before a_3
46x^3 - 6x^2 - 32x - 9 &= 0, & a_3 &= 2\\
x^3 - 94x^2 - 132x - 46 &= 0, & a_5 &= 95\\
3561x^3 - 9183x^2 - 191x - 1 &= 0, & a_6 &= 2,
\end{aligned}
\]
hence
\[
x = (1, 1, 3, 2, 1, 95, 2 \ldots)
\]
and the convergents:
%CONT>>
